\section{II. The resolved, calibrated heat-transition layer and the actual observed direction}\label{ii.-the-resolved-calibrated-heat-transition-layer-and-the-actual-observed-direction}

\subsection{II.1 What is inherited and what is being calculated}\label{ii.1-what-is-inherited-and-what-is-being-calculated}

The supplied H1--H15 retains the original full companion \(M=C_N+r_N\ell_N\), its original metric \(G_N\), the complete source minimum, and the uniform norm bound \(\|C_N\|_{G_N}\le B_hq\). It imports the profile asymptotic \[
 \log\epsilon_N^2=2q\psi(s_N)+O_h(k\log^2(q+2)),\quad
 \epsilon_N=\|r_N\ell_N\|_{G_N},\quad s_N=(N+1-q)/q.
\] Consequently it proves a heat step away from \(\sigma=\psi(s)\) at time \(e^{-2q\sigma}\), but it does not assign a value in that unresolved profile-centred layer. We do not claim to improve that profile error. The calculation below identifies the complete transition after centring at the \textbf{actual finite} \(\epsilon_N\), with a finite error valid before an asymptotic is applied.

\subsection{II.2 Finite identities and constants}\label{ii.2-finite-identities-and-constants}

Let \(M=C+r\ell\) on the given \(q\)-dimensional Hermitian space, and put \[
 \epsilon=\|r\ell\|_G>0,\quad \|C\|_G\le B,\quad
 d=B/\epsilon.
\] Assume a certified spectral-radius bound \(|\lambda(M)|\le R\) for the holomorphic trace. In the native instance this is the exact root bound \(R=k\sqrt{\delta^2+\gamma^2}\), including all primary multiplicities. No bound for the operator norm of \(M\) is substituted for it or inferred from it.

If \(s_1\ge\cdots\ge s_q\) are the singular values of \(M\), the rank-one minimum--maximum argument and the two triangle inequalities give \[
 |s_1-\epsilon|\le B,\qquad s_j\le B\quad(j\ge2).
 \tag{H1}
\] Set \(\tau=x/\epsilon^2\) for \(x\ge0\), \(y_j=s_j^2/\epsilon^2\). Then \[
 |y_1-1|\le2d+d^2,\qquad 0\le y_j\le d^2\ (j\ge2).
 \tag{H2}
\] For the full canonical heat gap, spectral logarithm, and resolvent response define \[
 \mathcal D(x)=\Re\operatorname{Tr}e^{-xM^2/\epsilon^2}
                  -\operatorname{Tr}e^{-xM^\dagger M/\epsilon^2},
\] \[
 \mathcal L(x)=\log\det(I+xM^\dagger M/\epsilon^2),\quad
 \mathcal A(x)=\operatorname{Tr}\frac{xM^\dagger M/\epsilon^2}
                                      {I+xM^\dagger M/\epsilon^2}.
\] Here the notation for the last quotient means multiplication by the inverse of the positive denominator, with the same operator in both factors. The full finite estimates are \[
 \boxed{|\mathcal D(x)-(1-e^{-x})|
 \le q\frac{xR^2}{\epsilon^2}e^{xR^2/\epsilon^2}
             +2xd+qx d^2,}
 \tag{H3}
\] \[
 \boxed{|\mathcal L(x)-\log(1+x)|\le2xd+qx d^2,\qquad
 |\mathcal A(x)-x/(1+x)|\le2xd+qx d^2.}
 \tag{H4}
\] \textbf{Proof.} The functions \(1-e^{-xy}\), \(\log(1+xy)\), and \(xy/(1+xy)\), as functions of \(y\ge0\), have derivatives in \([0,x]\). Apply H2 to the first singular value and use \(f(y)\le xy\) on the others. The holomorphic trace differs from \(q\) by at most \(q(xR^2/\epsilon^2)e^{xR^2/\epsilon^2}\), since on each full primary block the trace of the exponential is its multiplicity times the exponential at its eigenvalue. This proves H3--H4 without deleting a nilpotent block from the endomorphism.

Differentiation at the calibrated variable is also controlled: \[
 \boxed{|\mathcal D'(x)-e^{-x}|
 \le q\frac{R^2}{\epsilon^2}e^{xR^2/\epsilon^2}+2d+qd^2,}
 \tag{H5}
\] \[
 \boxed{|\mathcal L'(x)-(1+x)^{-1}|\le2d+qd^2.}
 \tag{H6}
\] For H5, \(y\mapsto y e^{-xy}\) has derivative \(e^{-xy}(1-xy)\) of absolute value at most one; its value on the remaining singular directions is at most \(y\). For H6, use the derivative \(1/(1+xy)^2\le1\) of \(y/(1+xy)\). The root-trace derivative is bounded separately as shown.

\subsection{II.3 Native consequences and what remains unknown at the profile threshold}\label{ii.3-native-consequences-and-what-remains-unknown-at-the-profile-threshold}

Under the supplied exponential \(\epsilon_N\) profile, \(qB_h^2q^2/\epsilon_N^2\) and \(qR_k^2/\epsilon_N^2\) tend to zero uniformly throughout the stated cutoff window. Thus H3--H6 give uniformly on compact \(x\)-intervals \[
 \mathcal D_N(x)\to1-e^{-x},\quad
 \mathcal L_N(x)\to\log(1+x),\quad
 \mathcal A_N(x)\to x/(1+x).
 \tag{H7}
\] The exact calibrated logarithmic centre is \[
 \widehat\psi_N=\frac{\log\epsilon_N^2}{2q}.
\] With \(x=e^z\) the time is \[
 \tau_N(z)=e^{-2q\widehat\psi_N+z},\qquad
 \sigma_N(z)=\widehat\psi_N-z/(2q).
\] The transition profile is \(1-e^{-e^z}\), the regularized log profile is \(\log(1+e^z)\), and the corresponding derivatives follow from H5--H6. Its intrinsic logarithmic width is \(1/q\). The supplied EIQ approximation locates its centre only to \(O_h(k\log^2q/q)\). Therefore this calculation does \textbf{not} assign a canonical value at the uncalibrated time \(e^{-2q\psi(s)}\).

For \(x=1\) the limiting full gap is \(1-e^{-1}\). Calibrating each of the four cutoffs separately gives four identical limiting constants, whose signed sum is zero; those are four \textbf{different physical heat times}. This is not a change to the source's common-time four-cutoff plateau. The latter must retain one common \(\tau\), and the exact different centres \(\widehat\psi_N\) remain in H3.

If \(\epsilon_N^2\) is known only in a positive certified interval, choose a declared reference \(E_0>0\) and set \(\tau=x/E_0\). Replace \(x\) in the monotone profiles by the interval \(x\epsilon_N^2/E_0\) and bound all H3 terms outward using the same interval. A profile centre estimated from a midpoint alone is not an exact native scale.

\subsection{II.4 The observed heat deficit measures a product of actual projection fractions}\label{ii.4-the-observed-heat-deficit-measures-a-product-of-actual-projection-fractions}

Let \(\Lambda:E\to B\) be the unchanged onto observation and \[
 Q_B=(\Lambda G^{-1}\Lambda^*)^{-1},\quad
 L_B=G^{-1}\Lambda^*Q_B,\quad P=L_B\Lambda.
\] Then \(L_B\) is an isometry and \(P\) is the \(G\)-orthogonal projection onto the attained observation complement. Let \(v=G^{-1}\ell^*\) and define \[
 \vartheta=
 \frac{\|Pr\|_G^2}{\|r\|_G^2}
 \frac{\|Pv\|_G^2}{\|v\|_G^2}\in[0,1].
 \tag{H8}
\] The exact observed operator is \[
 B_{\rm obs}=\Lambda ML_B=\Lambda CL_B+(\Lambda r)(\ell L_B).
\] Its bounded term has norm at most \(B\), and its rank-one term has squared norm \(\epsilon^2\vartheta\). This follows from \(\|\Lambda r\|_{Q_B}=\|Pr\|_G\) and the corresponding dual-row identity for \(v\). The projected rank-one term need not be square-zero.

For \(r_B=\dim B\ge1\), put \[
 \mathcal O(x)=r_B-\operatorname{Tr}
 e^{-x B_{\rm obs}^{\dagger_{Q_B}}B_{\rm obs}/\epsilon^2}.
\] The same singular-value calculation proves \[
 \boxed{|\mathcal O(x)-(1-e^{-x\vartheta})|
 \le2xd\sqrt\vartheta+r_Bx d^2
 \le2xd+r_Bxd^2.}
 \tag{H9}
\] It also proves \[
 \left|\log\det(I+xB_{\rm obs}^\dagger B_{\rm obs}/\epsilon^2)
 -\log(1+x\vartheta)\right|\le2xd+r_Bxd^2.
 \tag{H10}
\] The case \(r_B=0\) is empty and both observed quantities are zero. No holomorphic observed trace has been replaced by the original root trace: \(B_{\rm obs}\) is not generally an invariant restriction of \(M\).

Thus a limiting observed transition requires the actual product \(\vartheta_N\). If \(\vartheta_N\to\vartheta\), then \(\mathcal O_N(x)\to1-e^{-x\vartheta}\). Codimension alone does not evaluate this product. A one-dimensional kernel containing \(r\) or \(v\) makes \(\vartheta=0\), even if its codimension fraction vanishes; a kernel orthogonal to both leaves \(\vartheta=1\).

An interval for \(\mathcal O(x)\) gives an inverse angle certificate. If \([o_-,o_+]\) encloses it and \(e\) is the outward H9 error, set \[
 L=\max(0,o_--e),\qquad U=\min(1,o_++e).
\] When \(U<1\), \[
 \boxed{\vartheta\in
 \left[-\frac{\log(1-L)}x,-\frac{\log(1-U)}x\right]\cap[0,1],\qquad x>0.}
 \tag{H11}
\] If this interval fails to separate a needed value, no angle claim is made. It determines a product of two retained fractions, not their individual values or the complex current phases in the supplied C1--C4.

\subsection{II.5 An exact finite example for validation}\label{ii.5-an-exact-finite-example-for-validation}

Take \(G=I\), \(C=I_2\), and \(M=I_2+E e_1e_2^T\), \(E>0\). Its eigenvalues remain \((1,1)\) while \(\epsilon=E\). For projection onto the unit vector \((3/5,4/5)\), the observed operator is the scalar \(1+12E/25\), and \(\vartheta=144/625\). All heat and logarithmic enclosures in the checker are computed by rational power-series bounds, separately from H3--H10. This is a declared synthetic family; its entries are not substituted for native zeta moments.
